% Dokumentklasse
\documentclass[11pt,titlepage]{report} % Für längere Arbeiten mit Kapiteln und Abschnitten, geeignet für Berichte und Dokumentationen.

% Sprachunterstützung
\usepackage[ngerman]{babel} % Sorgt für sprachspezifische Regeln wie Silbentrennung und deutsche Bezeichnungen.

% Symbole und Icons
\usepackage{fontawesome} % Ermöglicht die Nutzung von Font Awesome Icons.

% Seitenlayout
\usepackage[a4paper,lmargin={3.5cm},rmargin={2.5cm},tmargin={3cm},bmargin={2cm}]{geometry} % Einstellung von Seitenrändern für A4-Papier.

% Mathematische Pakete
\usepackage{amssymb} % Enthält zusätzliche mathematische Symbole.
\usepackage{amsthm} % Unterstützung für Theoreme, Definitionen und Beweise.
\usepackage{amsmath} % Verbesserte Darstellung von mathematischen Formeln.

% Grafiken und Bilder
\usepackage{graphicx} % Ermöglicht das Einfügen von Bildern und Grafiken.

% Typografie
\usepackage{microtype} % Verbessert die Lesbarkeit und Silbentrennung.
\usepackage{fontspec} % Für die Verwendung von TrueType- und OpenType-Schriften.

% Hyperlinks
\usepackage[hidelinks]{hyperref}% Fügt klickbare Hyperlinks ein, ohne sie farblich hervorzuheben.

% Abkürzungen
\usepackage[footnote,printonlyused]{acronym} % Verwaltung von Abkürzungen mit Fußnotenoption.

% Kopf- und Fußzeilen
\usepackage{fancyhdr} % Ermöglicht die Gestaltung von Kopf- und Fußzeilen.

% Farben
\usepackage{xcolor} % Definiert und verwendet Farben im Dokument.

% Weitere mathematische Symbole
\usepackage{amsfonts} % Zusätzliche Schriftarten für mathematische Symbole.

% Tabellen
\usepackage{booktabs} % Verbesserte Gestaltung von Tabellen.
\usepackage{enumitem} % Anpassung von Listen und Aufzählungen.

% Literaturverwaltung
\usepackage[style=apa, backend=biber, style=authoryear, sorting=nyt, dashed=false, giveninits=true]{biblatex} % Für Literaturverzeichnisse nach APA-Stil.

% Zitate
\usepackage{csquotes} % Korrekte Darstellung von Zitaten, z. B. Anführungszeichen.

% Positionierung
\usepackage{float} % Präzise Platzierung von Tabellen und Abbildungen.

% Zeilenabstand
\usepackage[onehalfspacing]{setspace} % Erzeugt 1,5-fachen Zeilenabstand.

% Glossare
\usepackage[toc=false]{glossaries-extra} % Verwaltung von Glossaren ohne Inhaltsverzeichnis-Eintrag.

% Tabellenfarben
\usepackage{colortbl} % Ermöglicht das Einfärben von Tabellenzellen.

% Tabellen- und Spaltenformatierung
\usepackage{array} % Zusätzliche Optionen für Tabellen und Spalten.

% Diagramme und Visualisierungen
\usepackage{tikz} % Erstellung von Vektorgraphiken und Diagrammen.
\usepackage{pgfplots} % Erstellung von Plots und Grafiken.

\usetikzlibrary{shapes,arrows,positioning,fit,backgrounds,calc,matrix,decorations.pathreplacing,decorations.markings,shapes.geometric,mindmap}

% PDF-Integration
\usepackage{pdfpages} % Einfügen von ganzen PDF-Dokumenten.

% Querformat
\usepackage{pdflscape} % Ermöglicht die Nutzung von Seiten im Querformat.

% Nummerierungsanpassung
\usepackage{chngcntr} % Ermöglicht Änderungen der Nummerierung von Kapiteln und Abschnitten.

% Inhaltsverzeichnis
\usepackage{tocloft} % Anpassung des Layouts des Inhaltsverzeichnisses.

% Titelgestaltung
\usepackage{titlesec} % Anpassung von Kapitel- und Abschnittstiteln.

% Ausrichtung
\usepackage{ragged2e} % Ermöglicht den Blocksatz oder Flattersatz.

\pgfplotsset{compat=1.18}

% Zähler für Abbildungen und Tabellen an Abschnitte binden
\counterwithin{figure}{section}
\counterwithin{table}{section}
\counterwithin{equation}{section}

% Pfade für Bilder und andere Ressourcen
\graphicspath{{bilder/}}

% Arial als Hauptschrift (11pt)
\setmainfont{Arial}

% Bibliographie
\addbibresource{quellen.bib}

% Farben
\definecolor{link}{HTML}{0000EE}

% Absatzformatierung
\setlength{\parindent}{0pt}           % Keine Einrückung am Absatzbeginn

% Überschriftenabstände
% Anpassung \listoffigures
\renewcommand{\cftloftitlefont}{\Large\bfseries}
\renewcommand{\cftbeforeloftitleskip}{10pt}
\renewcommand{\cftafterloftitleskip}{10pt}

% Anpassung \listoftables
\renewcommand{\cftlottitlefont}{\Large\bfseries}
\renewcommand{\cftbeforelottitleskip}{10pt}
\renewcommand{\cftafterlottitleskip}{10pt}

% Anpassung \tableofcontents
\renewcommand{\cfttoctitlefont}{\Large\bfseries}
\renewcommand{\cftbeforetoctitleskip}{10pt}
\renewcommand{\cftaftertoctitleskip}{10pt}

% Abstand nach Nummern im Abbildungs- und Tabellenverzeichnis
\setlength{\cftfignumwidth}{3em}
\setlength{\cfttabnumwidth}{3em}

% Überschriftenformatierung
\titleformat{\chapter}
  {\fontsize{16}{19}\bfseries\raggedright}
  {\thechapter}{1em}{}
\titlespacing*{\chapter}{0pt}{0pt}{20pt}

\titleformat{\section}
  {\fontsize{14}{17}\bfseries\raggedright}
  {\thesection}{1em}{}

\titleformat{\subsection}
  {\fontsize{12}{14}\bfseries\raggedright}
  {\thesubsection}{1em}{}

\titleformat{\subsubsection}
  {\fontsize{12}{14}\bfseries\raggedright}
  {\thesubsubsection}{1em}{}

% Blocksatz für Fließtext
\justifying

% Bibliographie-Einstellungen
\DeclareNameAlias{sortname}{family-given}
\DeclareDelimFormat{nameyeardelim}{\addcomma\space}
\DeclareDelimFormat{multinamedelim}{\space\addcomma\space}
\DeclareDelimFormat{finalnamedelim}{\space\&\space}
\DefineBibliographyStrings{german}{andothers = {{et\,al\adddot}},}
\renewcommand{\mkbibnamefamily}[1]{\textsc{#1}}

% Benutzerdefinierte Zitatbefehle
\newcommand{\customcite}[2][]{%
  \ifblank{#1}{%
    \parencite{#2}%
  }{%
    \parencite[#1]{#2}%
  }%
}

% Fix lange URLs
\newcommand*\oldurlbreaks{}
\let\oldurlbreaks=\UrlBreaks
\renewcommand{\UrlBreaks}{\oldurlbreaks
  \do\a\do\b\do\c\do\d\do\e\do\f\do\g\do\h\do\i\do\j\do\k\do\l\do\m\do\n\do\o\do\p\do\q\do\r\do\s\do\t\do\u\do\v
  %\do\w%
  \do\x\do\y\do\z
  \do?\do&\do_
  \do\0\do\1\do\2\do\3\do\4\do\5\do\6\do\7\do\8\do\9
  \do:\do/
  }
\urlstyle{same}

% Befehle für Links und Verweise
\newcommand{\sct}[2]{\textcolor{link}{\hyperref[#1]{#2}}}
\newcommand{\gl}[1]{\gls{#1}}

% Inhaltsverzeichnis-Einrückung
\setlength{\cftsubsecindent}{1.5em}
\setlength{\cftsubsecindent}{3.0em}

% Glossar/Abkürzungsverzeichnis
\makenoidxglossaries
\newglossaryentry{ki}{
  name={KI},
  description={\emph{Künstliche Intelligenz} bezeichnet die Fähigkeit von Maschinen, menschenähnliche, intelligente Verhaltensweisen zu zeigen. Diese Technologie umfasst Methoden wie maschinelles Lernen, neuronale Netze und Expertensysteme.}
}
\glsaddall

% Customize page numbering on right hand side
\fancypagestyle{plain}{%
  \fancyhf{} %clear all header and footer fields
  \rhead{\thepage}
}%
\fancyhf{} %clear all header and footer fields

% Seitenstil für alle Seiten
\pagestyle{fancy}
\rhead{\thepage} % Seitenzahl oben rechts
\renewcommand{\headrulewidth}{0pt} % Keine Linie unter der Kopfzeile

% Anpassung der Verzeichnis-Überschriften
\renewcommand{\contentsname}{Inhaltsverzeichnis}
\renewcommand{\listfigurename}{Abbildungsverzeichnis}
\renewcommand{\listtablename}{Tabellenverzeichnis}

% Setze PDF-Metadaten
\hypersetup{
  pdftitle={Titel},
  pdfauthor={Martin Kribs},
  pdfsubject={Subjekt},
}

\begin{document}

% Titelseite ohne Seitenzahl
\begin{titlepage}
\thispagestyle{empty}

% Sperrvermerk
\chapter*{Sperrvermerk}
\thispagestyle{empty}

Diese Arbeit beinhaltet interne und vertrauliche Informationen der [Firmenname]. Die Weitergabe des Inhalts der Arbeit und eventuell beiliegender Zeichnungen, Daten und Informationen - auch auszugsweise - ist grundsätzlich untersagt. Ausnahmen bedürfen der schriftlichen Genehmigung der [Firmenname] und des Verfassers.

\vspace{2cm}

\noindent [Ort], den \today

\vspace{2cm}

\noindent\rule{6cm}{0.4pt}\\
Martin Kribs

\newpage

\begin{titlepage}
 
    \includegraphics[width=0.5\textwidth]{rh-logo}\\[2cm]
    
    \begin{minipage}{\textwidth}
    \hspace{3cm}\begin{minipage}{\dimexpr\textwidth-3cm\relax}
    {\Large\textbf{Fachbereich Ingenieurwesen}}\\
    
    {\Large Studiengang Informatik (B.Sc.)}\\
    \\
    \\
    
    
    {\Large\textbf{Subjekt}}\\
    
    \parbox{\textwidth}{\Large\raggedright Titel}\\
    \\
    \\
    
    {\normalsize\fontsize{11pt}{16.5pt}\selectfont vorgelegt von}\\
    
    {\Large Mustermensch}\\
    
    {\normalsize\fontsize{11pt}{16.5pt}\selectfont Matrikelnr. 111111111}\\
    \\
    
    {\normalsize Sommersemester 2025}
    \end{minipage}
    \end{minipage}
    
    \vfill
    
\end{titlepage}

\newpage
\newpage
\begin{titlepage}
    \centering
    
    {\huge Rheinische Hochschule Köln}\\[0.5cm]
    {\Large Fachbereich Ingenieurwesen}\\
    {\normalsize Studiengang: Informatik}\\[1cm]
    
    \includegraphics[width=0.3\textwidth]{rh-logo}\\[2cm]
    
    {\Large\textbf{Subjekt}}\\[1cm]
    
    {\Large Titel}\\[2cm]
    
    \normalsize
    Vorgelegt von: Martin Kribs \\
    Matrikelnummer: 653222007\\[2cm]
    
    1.Prüferin: Prüfer\\[0.3cm]
    
    Wintersemester 2024/25

    \vfill
    
\end{titlepage}

\end{titlepage}

\begin{onehalfspace}
    % Abkürzungsverzeichnis mit römischer Nummerierung
    \pagenumbering{Roman}
    \phantomsection
    \addcontentsline{toc}{section}{Abkürzungsverzeichnis}
    \label{sec:Abkuerzungsverzeichnis}
    \printnoidxglossary[title={Abkürzungsverzeichnis}, type=main, nonumberlist]
    \newpage
    
    % Inhaltsverzeichnis
    \phantomsection
    \addcontentsline{toc}{section}{Inhaltsverzeichnis}
    \tableofcontents
    \newpage
    
    % Tabellenverzeichnis
    \phantomsection
    \addcontentsline{toc}{section}{Tabellenverzeichnis}
    \listoftables
    \newpage
    
    % Abbildungsverzeichnis
    \phantomsection
    \addcontentsline{toc}{section}{Abbildungsverzeichnis}
    \listoffigures
    \newpage

    % Hauptteil-
    % Ab hier arabische Seitenzahlen
    \pagenumbering{arabic}

    % Kapitel einbinden
    \foreach \i in {10} {%
      \edef\FileName{abschnitte/\i-kapitel}%

      \IfFileExists{\FileName}{%
        \input{\FileName}
      }
      {%
        %file does not exist
      }
	  }

    % Literaturverzeichnis
    \newpage
    \phantomsection
    \addcontentsline{toc}{chapter}{Quellen- und Literaturverzeichnis}
    \label{chap:Quellen-und-Literaturverzeichnis}
    \printbibliography[title={Quellen- und Literaturverzeichnis}]
    
    % Anhänge
    \newpage
    \appendix
    \chapter{Anhang}
    \label{chap:Anhang}
    \section{Anhang}
\label{app:anhang}


    % Eigenständigkeitserklärung
    \newpage
    \chapter*{Eigenständigkeitserklärung}
\addcontentsline{toc}{chapter}{Eigenständigkeitserklärung}
Hiermit bestätige ich, dass ich die vorliegende Arbeit selbstständig verfasst und keine 
anderen Publikationen, Vorlagen und Hilfsmittel (z.B. künstliche Intelligenz) als die 
angegebenen benutzt habe. Alle Teile meiner Arbeit, die wortwörtlich oder dem Sinn nach 
anderen Werken entnommen sind, wurden unter Angabe der Quelle kenntlich gemacht. 
Gleiches gilt für von mir verwendete Internetquellen. Ich versichere, dass ich diese Arbeit 
oder nicht zitierte Teile daraus vorher nicht in einem anderen Prüfungsverfahren eingereicht 
habe. Mir ist bekannt, dass meine Arbeit zum Zwecke eines Plagiatsabgleichs mittels einer 
Plagiatserkennungssoftware auf eine ungekennzeichnete Übernahme von fremdem geistigen 
Eigentum sowie auf die Nutzung von künstlicher Intelligenz zur Texterstellung überprüft 
werden kann. Ich versichere, dass die elektronische Form meiner Arbeit mit der gedruckten 
Version identisch ist.


% \chapter*{Declaration of Originality}
\addcontentsline{toc}{chapter}{Declaration of Originality}
% I hereby confirm that I have independently written this work and have not used any 
% publications, templates, or aids (e.g. artificial intelligence) other than those I have 
% indicated. All parts of my work which have been taken literally or correspondingly from other
% publications have been duly acknowledged. This also applies to Internet sources. I confirm
% that I have not previously submitted this work or any unquoted parts thereof in any other 
% examination procedure. I am aware that my work may be checked for plagiarism by means 
% of plagiarism recognition software, as well as for the use of artificial intelligence for text
% creation, in order to verify the integrity of its written content. I also confirm that the 
% electronic form is identical to the printed version. \\

\vspace{2cm}

\begin{center}
\begin{tabular}{p{0.45\textwidth}p{0.45\textwidth}}
    [Ort], \today & \\
    \hline
    % Place, Date & Signature \\
    Ort, Datum & Unterschrift \\
\end{tabular}
\end{center}

\end{onehalfspace}
\end{document}
